\documentclass[10pt,a4paper]{article}
\usepackage[utf8]{inputenc}
\usepackage[top=2cm,bottom=2cm,left=2cm,right=2cm]{geometry}
\usepackage{amsmath}
\usepackage{amssymb}
\usepackage{array}
\usepackage[dvipsnames,svgnames]{xcolor}
\usepackage{hyperref}

% Anchor persamaan unik meski nomor di-reset per soal.
\newcounter{eqblock}
\newcommand{\reseteq}{%
  \stepcounter{eqblock}%
  \setcounter{equation}{0}%
  \renewcommand{\theHequation}{\arabic{eqblock}.\arabic{equation}}%
}

\newcommand{\dg}[1]{\textcolor{Red}{#1}}

\title{Tugas Minggu 11 : Matriks dan Operasinya}
\author{Matematika untuk Sistem Informasi}
\date{}

\begin{document}

\maketitle

\begin{enumerate}

%===========================================================
\item Diketahui matriks
\[
P = \begin{bmatrix} -2 & 0 & 1 \\ 2 & 1 & -3 \end{bmatrix} \quad \text{dan} \quad
Q = \begin{bmatrix} 0 & 3 & 1 \\ -1 & 1 & 2 \end{bmatrix}.
\]
Hitunglah:
\begin{enumerate}

%-----------------------------------------------------------
\item $P + Q$

\textbf{Jawab:}
\begin{itemize}
\item \textbf{Langkah 1: Periksa Ukuran}

$P$ dan $Q$ sama-sama berukuran $2 \times 3$, sehingga dapat dijumlahkan.

\item \textbf{Langkah 2: Jumlahkan Entri Seletak}

Hasilnya juga berukuran $2 \times 3$:
\reseteq
\begin{align}
P + Q &= \begin{bmatrix} -2 + 0 & 0 + 3 & 1 + 1 \\ 2 + (-1) & 1 + 1 & -3 + 2 \end{bmatrix}
= \begin{bmatrix} -2 & 3 & 2 \\ 1 & 2 & -1 \end{bmatrix}
\end{align}
% final: \begin{bmatrix} -2 & 3 & 2 \\ 1 & 2 & -1 \end{bmatrix}
\end{itemize}

%-----------------------------------------------------------
\item $-3Q$

\textbf{Jawab:}
\begin{itemize}
\item \textbf{Langkah 1: Kalikan Setiap Entri dengan $-3$}

Perkalian skalar tidak mengubah ukuran matriks:
\reseteq
\begin{align}
-3Q &= \begin{bmatrix} -3 \cdot 0 & -3 \cdot 3 & -3 \cdot 1 \\ -3 \cdot (-1) & -3 \cdot 1 & -3 \cdot 2 \end{bmatrix}
= \begin{bmatrix} 0 & -9 & -3 \\ 3 & -3 & -6 \end{bmatrix}
\end{align}
% final: \begin{bmatrix} 0 & -9 & -3 \\ 3 & -3 & -6 \end{bmatrix}
\end{itemize}

%-----------------------------------------------------------
\item $2Q - 3P$

\textbf{Jawab:}
\begin{itemize}
\item \textbf{Langkah 1: Hitung $2Q$ dan $3P$}

Kalikan setiap entri dengan skalarnya:
\reseteq
\begin{align}
2Q &= \begin{bmatrix} 0 & 6 & 2 \\ -2 & 2 & 4 \end{bmatrix}, \qquad
3P = \begin{bmatrix} -6 & 0 & 3 \\ 6 & 3 & -9 \end{bmatrix}
\end{align}

\item \textbf{Langkah 2: Kurangkan Entri Seletak}

Keduanya berukuran $2 \times 3$:
\reseteq
\begin{align}
2Q - 3P &= \begin{bmatrix} 0 - (-6) & 6 - 0 & 2 - 3 \\ -2 - 6 & 2 - 3 & 4 - (-9) \end{bmatrix}
= \begin{bmatrix} 6 & 6 & -1 \\ -8 & -1 & 13 \end{bmatrix}
\end{align}
% final: \begin{bmatrix} 6 & 6 & -1 \\ -8 & -1 & 13 \end{bmatrix}
\end{itemize}

\end{enumerate}

%===========================================================
\item Diketahui matriks-matriks berikut:
\[
A = \begin{bmatrix} 3 & 0 \\ -1 & 2 \\ 1 & 1 \end{bmatrix}, \quad
B = \begin{bmatrix} 1 & 4 & 2 \\ 3 & 1 & 5 \end{bmatrix}, \quad
C = \begin{bmatrix} 1 & 5 & 2 \\ -1 & 0 & 1 \\ 3 & 2 & 4 \end{bmatrix}, \quad
D = \begin{bmatrix} 6 & 1 & 3 \\ -1 & 1 & 2 \\ 4 & 1 & 3 \end{bmatrix}.
\]
Transpos yang sering dipakai:
\[
A^T = \begin{bmatrix} 3 & -1 & 1 \\ 0 & 2 & 1 \end{bmatrix}, \quad
B^T = \begin{bmatrix} 1 & 3 \\ 4 & 1 \\ 2 & 5 \end{bmatrix}, \quad
C^T = \begin{bmatrix} 1 & -1 & 3 \\ 5 & 0 & 2 \\ 2 & 1 & 4 \end{bmatrix}, \quad
D^T = \begin{bmatrix} 6 & -1 & 4 \\ 1 & 1 & 1 \\ 3 & 2 & 3 \end{bmatrix}.
\]
Hitunglah:
\begin{enumerate}

%-----------------------------------------------------------
\item $B^T A^T + 2D$

\textbf{Jawab:}
\begin{itemize}
\item \textbf{Langkah 1: Periksa Ukuran}

$B^T_{3 \times 2} \cdot A^T_{2 \times 3}$ menghasilkan matriks $3 \times 3$, sama dengan ukuran $2D$, sehingga keduanya dapat dijumlahkan.

\item \textbf{Langkah 2: Hitung $B^T A^T$ (Baris Kali Kolom)}

Setiap baris $B^T$ dikalikan dengan setiap kolom $A^T$:
\reseteq
\begin{align}
B^T A^T &= \begin{bmatrix} 1 & 3 \\ 4 & 1 \\ 2 & 5 \end{bmatrix} \begin{bmatrix} 3 & -1 & 1 \\ 0 & 2 & 1 \end{bmatrix}
= \begin{bmatrix} 3 + 0 & -1 + 6 & 1 + 3 \\ 12 + 0 & -4 + 2 & 4 + 1 \\ 6 + 0 & -2 + 10 & 2 + 5 \end{bmatrix}
= \begin{bmatrix} 3 & 5 & 4 \\ 12 & -2 & 5 \\ 6 & 8 & 7 \end{bmatrix}
\end{align}

\item \textbf{Langkah 3: Tambahkan $2D$}

Dengan $2D = \begin{bmatrix} 12 & 2 & 6 \\ -2 & 2 & 4 \\ 8 & 2 & 6 \end{bmatrix}$:
\reseteq
\begin{align}
B^T A^T + 2D &= \begin{bmatrix} 3 + 12 & 5 + 2 & 4 + 6 \\ 12 - 2 & -2 + 2 & 5 + 4 \\ 6 + 8 & 8 + 2 & 7 + 6 \end{bmatrix}
= \begin{bmatrix} 15 & 7 & 10 \\ 10 & 0 & 9 \\ 14 & 10 & 13 \end{bmatrix}
\end{align}

\item \textbf{Langkah 4: Cek dengan Sifat $(AB)^T = B^T A^T$}

Hitung $AB$ lalu transposkan:
\reseteq
\begin{align}
AB &= \begin{bmatrix} 3 & 12 & 6 \\ -1 + 6 & -4 + 2 & -2 + 10 \\ 1 + 3 & 4 + 1 & 2 + 5 \end{bmatrix}
= \begin{bmatrix} 3 & 12 & 6 \\ 5 & -2 & 8 \\ 4 & 5 & 7 \end{bmatrix} \\
(AB)^T &= \begin{bmatrix} 3 & 5 & 4 \\ 12 & -2 & 5 \\ 6 & 8 & 7 \end{bmatrix} = B^T A^T
\end{align}
% final: \begin{bmatrix} 15 & 7 & 10 \\ 10 & 0 & 9 \\ 14 & 10 & 13 \end{bmatrix}
\end{itemize}

%-----------------------------------------------------------
\item $\operatorname{tr}(4D^T - C)$

\textbf{Jawab:}
\begin{itemize}
\item \textbf{Langkah 1: Hitung $4D^T - C$}

Kalikan $D^T$ dengan 4, lalu kurangkan $C$ entri demi entri:
\reseteq
\begin{align}
4D^T - C &= \begin{bmatrix} 24 & -4 & 16 \\ 4 & 4 & 4 \\ 12 & 8 & 12 \end{bmatrix} - \begin{bmatrix} 1 & 5 & 2 \\ -1 & 0 & 1 \\ 3 & 2 & 4 \end{bmatrix}
= \begin{bmatrix} \dg{23} & -9 & 14 \\ 5 & \dg{4} & 3 \\ 9 & 6 & \dg{8} \end{bmatrix}
\end{align}

\item \textbf{Langkah 2: Jumlahkan Diagonal Utama}

Entri diagonal ditandai merah:
\reseteq
\begin{align}
\operatorname{tr}(4D^T - C) &= 23 + 4 + 8 = 35
\end{align}

\item \textbf{Langkah 3: Cek dengan Sifat Trace}

Karena $\operatorname{tr}(kA) = k\operatorname{tr}(A)$, $\operatorname{tr}(A^T) = \operatorname{tr}(A)$, dan trace bersifat linear:
\reseteq
\begin{align}
\operatorname{tr}(4D^T - C) &= 4\operatorname{tr}(D) - \operatorname{tr}(C) = 4(6 + 1 + 3) - (1 + 0 + 4) = 40 - 5 = 35
\end{align}
% final: 35
\end{itemize}

%-----------------------------------------------------------
\item $(2C^T - D)A$

\textbf{Jawab:}
\begin{itemize}
\item \textbf{Langkah 1: Periksa Ukuran}

$(2C^T - D)$ berukuran $3 \times 3$ dan $A$ berukuran $3 \times 2$, sehingga hasilnya berukuran $3 \times 2$.

\item \textbf{Langkah 2: Hitung $2C^T - D$}

Kalikan $C^T$ dengan 2, lalu kurangkan $D$:
\reseteq
\begin{align}
2C^T - D &= \begin{bmatrix} 2 & -2 & 6 \\ 10 & 0 & 4 \\ 4 & 2 & 8 \end{bmatrix} - \begin{bmatrix} 6 & 1 & 3 \\ -1 & 1 & 2 \\ 4 & 1 & 3 \end{bmatrix}
= \begin{bmatrix} -4 & -3 & 3 \\ 11 & -1 & 2 \\ 0 & 1 & 5 \end{bmatrix}
\end{align}

\item \textbf{Langkah 3: Kalikan dengan $A$}

Setiap baris $(2C^T - D)$ dikalikan dengan kedua kolom $A$:
\reseteq
\begin{align}
(2C^T - D)A &= \begin{bmatrix} -4 & -3 & 3 \\ 11 & -1 & 2 \\ 0 & 1 & 5 \end{bmatrix} \begin{bmatrix} 3 & 0 \\ -1 & 2 \\ 1 & 1 \end{bmatrix}
= \begin{bmatrix} -12 + 3 + 3 & 0 - 6 + 3 \\ 33 + 1 + 2 & 0 - 2 + 2 \\ 0 - 1 + 5 & 0 + 2 + 5 \end{bmatrix}
= \begin{bmatrix} -6 & -3 \\ 36 & 0 \\ 4 & 7 \end{bmatrix}
\end{align}
% final: \begin{bmatrix} -6 & -3 \\ 36 & 0 \\ 4 & 7 \end{bmatrix}
\end{itemize}

%-----------------------------------------------------------
\item $(2C^T - 3D^T)^T$

\textbf{Jawab:}
\begin{itemize}
\item \textbf{Langkah 1: Hitung $2C^T - 3D^T$}

Kalikan $C^T$ dengan 2 dan $D^T$ dengan 3, lalu kurangkan:
\reseteq
\begin{align}
2C^T - 3D^T &= \begin{bmatrix} 2 & -2 & 6 \\ 10 & 0 & 4 \\ 4 & 2 & 8 \end{bmatrix} - \begin{bmatrix} 18 & -3 & 12 \\ 3 & 3 & 3 \\ 9 & 6 & 9 \end{bmatrix}
= \begin{bmatrix} -16 & 1 & -6 \\ 7 & -3 & 1 \\ -5 & -4 & -1 \end{bmatrix}
\end{align}

\item \textbf{Langkah 2: Transposkan}

Baris menjadi kolom:
\reseteq
\begin{align}
(2C^T - 3D^T)^T &= \begin{bmatrix} -16 & 7 & -5 \\ 1 & -3 & -4 \\ -6 & 1 & -1 \end{bmatrix}
\end{align}

\item \textbf{Langkah 3: Cara Cepat dengan Sifat Transpos}

Dari $(A - B)^T = A^T - B^T$, $(kA)^T = kA^T$, dan $(A^T)^T = A$ diperoleh $(2C^T - 3D^T)^T = 2C - 3D$:
\reseteq
\begin{align}
2C - 3D &= \begin{bmatrix} 2 - 18 & 10 - 3 & 4 - 9 \\ -2 + 3 & 0 - 3 & 2 - 6 \\ 6 - 12 & 4 - 3 & 8 - 9 \end{bmatrix}
= \begin{bmatrix} -16 & 7 & -5 \\ 1 & -3 & -4 \\ -6 & 1 & -1 \end{bmatrix}
\end{align}
Hasilnya sama, tanpa perlu menghitung transpos sama sekali.
% final: \begin{bmatrix} -16 & 7 & -5 \\ 1 & -3 & -4 \\ -6 & 1 & -1 \end{bmatrix}
\end{itemize}

%-----------------------------------------------------------
\item $\operatorname{tr}(B^T A^T + 2C^T)$

\textbf{Jawab:}
\begin{itemize}
\item \textbf{Langkah 1: Gunakan $B^T A^T$ dari (a)}

Tambahkan $2C^T$:
\reseteq
\begin{align}
B^T A^T + 2C^T &= \begin{bmatrix} 3 & 5 & 4 \\ 12 & -2 & 5 \\ 6 & 8 & 7 \end{bmatrix} + \begin{bmatrix} 2 & -2 & 6 \\ 10 & 0 & 4 \\ 4 & 2 & 8 \end{bmatrix}
= \begin{bmatrix} \dg{5} & 3 & 10 \\ 22 & \dg{-2} & 9 \\ 10 & 10 & \dg{15} \end{bmatrix}
\end{align}

\item \textbf{Langkah 2: Jumlahkan Diagonal Utama}

Entri diagonal ditandai merah:
\reseteq
\begin{align}
\operatorname{tr}(B^T A^T + 2C^T) &= 5 + (-2) + 15 = 18
\end{align}

\item \textbf{Langkah 3: Cek dengan Sifat Trace}

Karena $B^T A^T = (AB)^T$ dan transpos tidak mengubah trace:
\reseteq
\begin{align}
\operatorname{tr}(B^T A^T + 2C^T) &= \operatorname{tr}(AB) + 2\operatorname{tr}(C) = (3 - 2 + 7) + 2(1 + 0 + 4) = 8 + 10 = 18
\end{align}
% final: 18
\end{itemize}

%-----------------------------------------------------------
\item $C^T D^T - (DC)^T$

\textbf{Jawab:}
\begin{itemize}
\item \textbf{Langkah 1: Hitung $C^T D^T$}

Setiap baris $C^T$ dikalikan dengan setiap kolom $D^T$:
\reseteq
\begin{align}
C^T D^T &= \begin{bmatrix} 1 & -1 & 3 \\ 5 & 0 & 2 \\ 2 & 1 & 4 \end{bmatrix} \begin{bmatrix} 6 & -1 & 4 \\ 1 & 1 & 1 \\ 3 & 2 & 3 \end{bmatrix} \nonumber \\
&= \begin{bmatrix} 6 - 1 + 9 & -1 - 1 + 6 & 4 - 1 + 9 \\ 30 + 0 + 6 & -5 + 0 + 4 & 20 + 0 + 6 \\ 12 + 1 + 12 & -2 + 1 + 8 & 8 + 1 + 12 \end{bmatrix}
= \begin{bmatrix} 14 & 4 & 12 \\ 36 & -1 & 26 \\ 25 & 7 & 21 \end{bmatrix}
\end{align}

\item \textbf{Langkah 2: Hitung $DC$ dan $(DC)^T$}

Setiap baris $D$ dikalikan dengan setiap kolom $C$, lalu hasilnya ditransposkan:
\reseteq
\begin{align}
DC &= \begin{bmatrix} 6 & 1 & 3 \\ -1 & 1 & 2 \\ 4 & 1 & 3 \end{bmatrix} \begin{bmatrix} 1 & 5 & 2 \\ -1 & 0 & 1 \\ 3 & 2 & 4 \end{bmatrix} \nonumber \\
&= \begin{bmatrix} 6 - 1 + 9 & 30 + 0 + 6 & 12 + 1 + 12 \\ -1 - 1 + 6 & -5 + 0 + 4 & -2 + 1 + 8 \\ 4 - 1 + 9 & 20 + 0 + 6 & 8 + 1 + 12 \end{bmatrix}
= \begin{bmatrix} 14 & 36 & 25 \\ 4 & -1 & 7 \\ 12 & 26 & 21 \end{bmatrix} \\
(DC)^T &= \begin{bmatrix} 14 & 4 & 12 \\ 36 & -1 & 26 \\ 25 & 7 & 21 \end{bmatrix}
\end{align}

\item \textbf{Langkah 3: Kurangkan}

Kedua matriks identik, sehingga
\reseteq
\begin{align}
C^T D^T - (DC)^T &= \begin{bmatrix} 0 & 0 & 0 \\ 0 & 0 & 0 \\ 0 & 0 & 0 \end{bmatrix} = O_{3 \times 3}
\end{align}

\item \textbf{Langkah 4: Penjelasan dengan Sifat Transpos}

Menurut sifat $(AB)^T = B^T A^T$, berlaku $(DC)^T = C^T D^T$. Jadi $C^T D^T - (DC)^T$ \emph{selalu} matriks nol, untuk sembarang $C$ dan $D$ yang dapat dikalikan.
% final: \begin{bmatrix} 0 & 0 & 0 \\ 0 & 0 & 0 \\ 0 & 0 & 0 \end{bmatrix}
\end{itemize}

\end{enumerate}

%===========================================================
\item Tentukan nilai $a$, $b$, $c$, dan $d$ pada persamaan matriks berikut:
\[
\begin{bmatrix} a - b & b + c \\ 3d + c & 2a - 4d \end{bmatrix} = \begin{bmatrix} 8 & 1 \\ 7 & 6 \end{bmatrix}.
\]

\textbf{Jawab:}
\begin{itemize}
\item \textbf{Langkah 1: Samakan Entri Seletak}

Kedua matriks berukuran $2 \times 2$, sehingga diperoleh empat persamaan:
\reseteq
\begin{align}
a - b &= 8 \\
b + c &= 1 \\
3d + c &= 7 \\
2a - 4d &= 6
\end{align}

\item \textbf{Langkah 2: Eliminasi $a$ dari (1) dan (4)}

Kalikan (1) dengan 2 menjadi $2a - 2b = 16$, lalu kurangkan dari (4):
\begin{align}
(2a - 4d) - (2a - 2b) &= 6 - 16 \;\Rightarrow\; 2b - 4d = -10 \nonumber \\
b - 2d &= -5
\end{align}

\item \textbf{Langkah 3: Eliminasi $c$ dari (2) dan (3)}

Kurangkan (2) dari (3):
\begin{align}
(3d + c) - (b + c) &= 7 - 1 \;\Rightarrow\; 3d - b = 6 \nonumber \\
b - 3d &= -6
\end{align}

\item \textbf{Langkah 4: Eliminasi $b$ dari (5) dan (6), lalu Substitusi Balik}

Kurangkan (6) dari (5):
\begin{align}
(b - 2d) - (b - 3d) &= -5 - (-6) \;\Rightarrow\; d = 1 \nonumber \\
b &= -5 + 2d = -3 \quad \text{(Dari (5))} \nonumber \\
a &= 8 + b = 5 \quad \text{(Dari (1))} \nonumber \\
c &= 1 - b = 4 \quad \text{(Dari (2))} \nonumber
\end{align}

\item \textbf{Langkah 5: Verifikasi}

$a - b = 5 + 3 = 8$, $b + c = -3 + 4 = 1$, $3d + c = 3 + 4 = 7$, $2a - 4d = 10 - 4 = 6$. Semua sesuai (soal ini sama dengan Latihan Mandiri~8.2 pada Topik~9).

Jadi $a = 5$, $b = -3$, $c = 4$, dan $d = 1$.
% final: a = 5
% final: b = -3
% final: c = 4
% final: d = 1
\end{itemize}

\end{enumerate}

\end{document}
